\documentclass[10pt,a4paper]{amsart}
\usepackage[margin=19mm]{geometry}
\usepackage{amsmath,amssymb,mathtools}
\usepackage[scr=rsfs,cal=euler]{mathalfa}
\usepackage{xfrac}
\usepackage[unicode,hidelinks,bookmarks=false]{hyperref}
\newcommand{\ie}{i.e.\ }
\newcommand{\cf}{cf.\ }
\newcommand{\resp}{respectively\ }
\newcommand{\Subsection}{\subsection}
\providecommand{\nobreakdash}{\nobreak\hbox{-}}
\setlength{\emergencystretch}{2em}
\multlinegap0pt
\usepackage{amssymb}
\usepackage{enumerate}
\usepackage{mathtools}
\usepackage{float}
\newtheorem{theorem}{Theorem}
\title{On Modular Invariants of Truncated Polynomial Rings}
\author{Hoang Le Xuan}
\date{Source: arXiv:2605.30397v1 (CC BY 4.0)}
\begin{document}
\maketitle

\noindent\emph{Source and licence.} On Modular Invariants of Truncated Polynomial Rings. Version arXiv:2605.30397v1 (CC BY 4.0), distributed by arXiv under the Creative Commons Attribution 4.0 International licence, http://creativecommons.org/licenses/by/4.0/, as recorded in the arXiv licence metadata for this exact version and shown on the abstract page https://arxiv.org/abs/2605.30397v1. Full licence notice: \texttt{licenses/arxiv-2605.30397-CC-BY-4.0.txt}.\par\smallskip

\noindent\textit{Editorial scope.} Counted unit: the article's theorem main\_rewrite (source0.tex lines 906--912). The article's complete original proof of it is delivered with it (source0.tex lines 913--941, one \texttt{proof} environment of 29 lines). No mathematics is written, repaired or re-derived: the statement and the proof are the article's own lines, in the article's own notation, and no retained line is substituted or re-flowed.  Retained uncounted as the source-local context the proof needs: lines 862--873.  Nothing was omitted from any retained span, so the package carries no omission record.  No retained line contains a citation, so the wrapper carries no bibliography and no bibliography entry.  No retained line contains a figure environment, an image-inclusion command or drawing code, so no image is delivered and the package declares no asset dependency.  Editorial formatting only, in addition: an AMS \texttt{amsart} preamble that restates the article's own declarations and macros used by the retained lines, namely the environments \texttt{theorem} on separate counters, together with the standard packages those lines need.  The retained environments print in this document as Theorem 1 in order of appearance, whereas the article prints the same items with the numbering of its own full document.  The complete native source of the article as distributed in the arXiv e-print for this exact version is bundled unchanged as \texttt{sources/201670/source0.tex}.
\begin{theorem}\label{gen_of_7.25}
Let $U$ be a vector subspace of $V$, and let $\mu\subseteq \lambda$ be partitions.  If $\lambda/\mu$ has strictly less than $\dim(U)$ nonzero rows, then we have
\begin{equation*}
    S_{\lambda/\mu}(V) = \sum\limits_{\substack{L\subseteq U\\
    \dim(L) = 1}}S_{\lambda/\mu}(V/L).  
\end{equation*}
Similarly, if $\lambda/\mu$ has strictly less than $\dim (U)$ nonzero columns, then we have 
\begin{equation*}
    T_{\lambda/\mu}(V) = \sum\limits_{\substack{L\subseteq U\\
    \dim(L) = 1}}T_{\lambda/\mu}(V/L). 
\end{equation*}
\end{theorem}

\begin{theorem}\label{main_rewrite}
Let $V$ be an $n$-dimensional vector space over $\mathbb{F}_q$, and let $\mathcal{F}$ denote the the set of all complete flags\index{complete flag} $\underline{V} = (V=V_0 \supset \ldots \supset V_{n-1} \supset V_n = 0$). For any partition $\lambda$ of length at most $n$, we have the following formula
\begin{equation*}  
    S_\lambda(V) = (-1)^{|\lambda|} \sum\limits_{\underline{V} \in \mathcal{F}}\prod\limits_{i=1}^{n} \pi_i (\underline{V})^{\frac{q^{\lambda_{i}}-1}{q-1}},
\end{equation*}
where for a subspace $W \subset \mathbb{F}_q[V]$, $\pi(W)$ is the product of all non-zero vectors in $W$, and for a full flag $\underline{V} = (0 = V_n \subsetneq V_{n-1}\subsetneq \cdots \subsetneq V_1\subsetneq V_0 = V)$, and $\pi_i(\underline{V}) = \pi(V_{i-1}/V_i)$ for $1 \leq i \leq n$.
\end{theorem}

\begin{proof}
If $\ell (\lambda) = n$, in other words $\lambda_n > 0$, then $\lambda_1' = n$ and $E_{\lambda_1' - 1+j} = 0$ unless $j=1$, in which case it is $E_n$. Thus 
\begin{align*}
S_{\lambda}(V) =& \det (\varphi^{j-1} E_{\lambda_i' - i +j})\\
=&  E_n (V) \det (\varphi^{j-1} E_{\lambda_i' - i +j})_{i,j \geq 2}\\
=& E_n (V) \varphi S_{(\lambda_1-1, \ldots, \lambda_n-1)} (V). 
\end{align*}
Repeating this process, we get 
\begin{equation*}
S_{\lambda} (V) = E_n(V)^{\frac{q^{\lambda_n}-1}{q-1}} \varphi^{\lambda_n} S_{(\lambda_1 - \lambda_n, \ldots, \lambda_{n-1}-\lambda_n, 0)} (V). 
\end{equation*}
Now the length of $(\lambda_1 - \lambda_n, \ldots, \lambda_{n-1} - \lambda_n,0)$ is strictly less than $n = \dim V$. We can apply Theorem \ref{gen_of_7.25} in the case $\mu = 0$: 
\begin{equation*}
S_{\lambda} (V) =  E_n(V)^{\frac{q^{\lambda_n}-1}{q-1}} \sum_{V_{n-1} \leq V, \dim V_{n-1} =1} S_{(\lambda_1 - \lambda_n, \ldots, \lambda_{n-1} - \lambda_n)}^{q^{\lambda_n}} (V/V_{n-1}).
\end{equation*}
Continue this process for each $(n-1)$-dimensional spaces $V/V_{n-1}$ and partitions $(\lambda_1- \lambda_n, \ldots, \lambda_{n-1}-\lambda_n)$ of length at most $(n-1)$, we can write $S_{\lambda}(V)$ as a sum over all complete flags: 
\begin{equation*}
\sum_{\underline{V} \in \mathcal{F}} E_{n}(V)^{\frac{q^{\lambda_n}-1}{q-1}}   E_{n-1}(V/V_{n-1})^{\frac{q^{\lambda_{n-1}} -q^{\lambda_n}}{q-1}}  E_{n-2}(V/V_{n-2})^{\frac{q^{\lambda_{n-2}} -q^{\lambda_{n-1}}}{q-1}}  \ldots E_{1}(V/V_{1})^{\frac{q^{\lambda_{1}}-q^{\lambda_2}}{q-1}}   
\end{equation*}
The result now follows from the fact that
\begin{equation*} 
E_{\dim V/U} (V/U) = (-1)^{\dim (V/U)} \pi (V/U) =(-1)^{\dim V - \dim U} \frac{\pi (V)}{\pi (U)}. 
\end{equation*}
For the general case, if $\lambda/\mu$ has exactly $n$ rows, then $\lambda_i - \mu_i > 0$ for all $1 \leq i \leq n$. It follows that $\lambda_1' = n$. If $\mu_1' = n$ then the first column of $S_{\lambda/\mu} (V)$ consists of all zero except in entries $(1,1)$ where it equals $\varphi^{-n} E_0 = 1$. We then have 
\begin{equation*}
S_{\lambda/\mu} (V) = S_{(\lambda_1-1, \ldots, \lambda_n-1)/(\mu_1-1, \ldots, \mu_n-1)} (V). 
\end{equation*}
Continue this process $\mu_n$ times, we can assume that $\mu_1' <n$ (or equivalently, $\mu_n = 0$).   
\end{proof}

\end{document}
